\chapter{The Agentic SDK Landscape: Frameworks \& Architectures}

\begin{quote}
\textit{``In the early days of LLMs, we wrote raw prompt-response loops and hoped for the best. Today, enterprise agents are built on standardized software development kits (SDKs) that manage conversation state, orchestrate multi-agent collaboration, enforce deterministic execution, and persist memory across service crashes. Choosing the right SDK is not a matter of developer preference --- it is a choice of system architecture.''}
\end{quote}

\begin{center}
\noindent\rule{0.8\textwidth}{0.4pt} \\
\vspace{0.3em}
\textbf{Companion Masterclass:} This chapter is mapped directly to the \textit{AI Agent Masterclass} repository. Visit the repository to access hands-on Jupyter Notebooks for all covered frameworks: \\
\url{https://github.com/umang-algo/AI-Agent-Masterclass} \\
\noindent\rule{0.8\textwidth}{0.4pt}
\end{center}

When building autonomous agents, developers rarely write raw API loops from scratch. Instead, they leverage standardized SDKs to structure agent roles, state memory, and transitions.

\section{The Five SDK Ecosystems}

As analyzed in the AI Agent Masterclass, the modern landscape is divided into five dominant ecosystems:

\subsection{1. The Microsoft Ecosystem: AutoGen \& Semantic Kernel}
Microsoft's SDK landscape is optimized for **Enterprise Orchestration** and **Message-Passing Multi-Agent Collaboration**.
\begin{itemize}
    \item \textbf{AutoGen (The Collaborative Swarm):} Focuses on the Actor Model of computing. Specialized agents are defined with specific system instructions and communicate via message-passing. A \texttt{GroupChatManager} acts as an arbiter, utilizing an LLM or deterministic rules to select the next speaker, allowing complex debates (e.g., CEO, Risk Manager, and Analyst) to run autonomously.
    \item \textbf{Semantic Kernel:} Focuses on integration. It models LLM capabilities as plugins (native functions or semantic prompts) that can be chained together using planners, making it the choice for integrating with enterprise C\# and Python services.
\end{itemize}

\subsection{2. The LangChain Ecosystem: LangChain \& LangGraph}
The LangChain ecosystem is the industry standard for stateful, cyclic graphs and customizable retrieval pipelines.
\begin{itemize}
    \item \textbf{LangChain (Knowledge Grounding/RAG):} Provides standard abstractions for documents, text splitters, vector stores, and retriever interfaces to build retrieval pipelines.
    \item \textbf{LangGraph (Stateful Graph Orchestration):} Models agents as stateful graphs where nodes are Python functions (agents/tools) and edges define transitions. Unlike linear chains, LangGraph allows cycles (loops), making it perfect for ReAct-style self-healing processes (Chapter 3). It features native state checkpointing for human-in-the-loop approvals.
\end{itemize}

\subsection{3. Native Provider Tooling: OpenAI \& Google Gen AI SDK}
Model providers offer optimized native SDKs that bypass third-party framework overhead.
\begin{itemize}
    \item \textbf{OpenAI Assistants API:} A server-side agent engine. OpenAI hosts the thread history, vector indexes (File Search), and code sandbox, reducing client-side code complexity.
    \item \textbf{Google Gen AI SDK (Agent Development Kit):} Built natively for Gemini 2.0 models. It provides features like **Search Grounding** (where the model verifies facts against live Google Search indices before responding) and native function-calling with low latency.
\end{itemize}

\subsection{4. Specialized Open-Source: CrewAI \& HuggingFace smolagents}
Specialized open-source frameworks focus on usability, roleplay modeling, and lightweight execution.
\begin{itemize}
    \item \textbf{CrewAI (Role-Playing Collaboration):} Models systems around human organizational metaphors: Crews, Tasks, and Agents. Each agent is given a specific role, backstory, and target goal. It manages transitions (sequential or hierarchical) to build multi-role pipelines.
    \item \textbf{HuggingFace smolagents (Code-Native Agents):} A shift away from JSON tool-calling. In smolagents, the agent outputs executable Python code directly. The local runtime runs this code inside a secure environment to call tools, reducing parsing errors.
\end{itemize}

\subsection{5. Enterprise \& Low-Code: Dify \& n8n}
For visual development and rapid enterprise integration, low-code platforms decouple logic from code.
\begin{itemize}
    \item \textbf{Dify (Headless Agent Engine):} A visual canvas to construct agent flows, prompt templates, and RAG pipelines. Once designed, the agent is exposed via a headless REST API.
    \item \textbf{n8n (Autonomous Action Layer):} An workflow automation tool with extensive node integrations. n8n allows agents to trigger webhooks, update CRMs, or route alerts through interactive visual nodes.
\end{itemize}

\section{Ecosystem Comparison Matrix}

The table below summarizes the trade-offs across these agent development ecosystems:

\begin{table}[h!]
\centering
\small
\begin{tabular}{|l|p{3.2cm}|p{3.2cm}|c|c|}
\hline
\textbf{Framework} & \textbf{State Persistence} & \textbf{Graph Topology} & \textbf{Parallel Execution} & \textbf{Primary Focus} \\ \hline
\textbf{LangGraph} & Database Checkpoints & Stateful Cyclic Graph & Yes & Control \& Reliability \\ \hline
\textbf{AutoGen} & Conversational DB Logs & Dynamic Conversation & Yes & Collaborative Swarms \\ \hline
\textbf{CrewAI} & Shared Context Class & Sequential / Hierarchical & No & Usability \& Persona \\ \hline
\textbf{smolagents} & Code-based State & Single loop / Custom & No & Code-Native Execution \\ \hline
\textbf{OpenAI API} & Managed Server Threads & Linear Run Loops & No & Managed Services \\ \hline
\textbf{Google ADK} & Gemini Grounding Store & Declarative / Tool List & Yes & Search Grounding \\ \hline
\textbf{Dify / n8n} & Visual Platform State & Stateful Workflow DAG & Yes & Integration \& Visual \\ \hline
\end{tabular}
\caption{Comparison of Agentic SDK Architectures.}
\end{table}

\section{Stateful Agent Graph Architecture}

The diagram below represents the unified architecture of a stateful agent graph orchestrator (e.g., LangGraph / AutoGen hybrid model). State is maintained centrally, and transitions are governed by nodes and edge routers.

\begin{figure}[ht]
\centering
\begin{tikzpicture}[
    scale=0.85, transform shape,
    node distance=0.6cm and 1.2cm,
    layer/.style={draw, rectangle, rounded corners, minimum height=0.75cm, minimum width=2.8cm, align=center, fill=blue!5, font=\small},
    dbbox/.style={draw, cylinder, shape border rotate=90, minimum height=0.8cm, minimum width=1.5cm, align=center, fill=orange!8, font=\small},
    router/.style={draw, shape=diamond, aspect=2.2, align=center, fill=yellow!15, font=\small},
    arrow/.style={thick,->,>=stealth}
]

% Column 1: Input, State, Checkpointer
\node[layer, fill=gray!15] (input) {User Query / Trigger};
\node[layer, below=of input] (state) {Central State \\ (Shared Thread State)};
\node[dbbox, below=of state] (checkpointer) {State Checkpointer \\ (PostgreSQL / Redis)};

% Column 2: Agent Nodes
\node[layer, right=1.6cm of state, yshift=1.2cm] (nodeA) {Agent Node A \\ (e.g., Coder Agent)};
\node[layer, right=1.6cm of state, yshift=-1.2cm] (nodeB) {Agent Node B \\ (e.g., Tester Agent)};

% Column 3: Router & Sandbox (Aligned based on Midpoint of Column 2)
\coordinate (mid) at ($(nodeA.east)!0.5!(nodeB.east)$);
\node[router, right=1.6cm of mid] (routerA) {Routing Decision \\ (LLM / Assert)};
\node[layer, below=of routerA, yshift=-0.4cm] (tool) {Tool Execution \\ Sandbox};

% Connections
\draw[arrow] (input) -- (state);
\draw[arrow, <->, dashed] (state) -- (checkpointer);
\draw[arrow] (state) -- (nodeA);
\draw[arrow] (state) -- (nodeB);
\draw[arrow] (nodeA) -- (routerA);
\draw[arrow] (nodeB) -- (routerA);
\draw[arrow] (routerA) -- node[right, font=\tiny] {use\_tool} (tool);

% Feedback Loops (Routed around the grid to save space)
\draw[arrow, dashed] (routerA.north) -- ++(0,1.2) -| node[above, xshift=2.5cm, font=\tiny] {Loop / Update State} (state.120);
\draw[arrow, dashed] (tool.south) -- ++(0,-1.0) -| node[below, xshift=2.5cm, font=\tiny] {Update State} (state.240);

\end{tikzpicture}
\caption{Stateful Agent Graph Orchestration Loop. The Checkpointer serializes graph states dynamically at every transition, enabling transactional safety and crash recovery.}
\end{figure}

\section{Asynchronous Graph State Serialization}

For industrial deployment, agents must survive application crashes. Runtimes resolve this by serializing the execution state graph into a database checkpoint before every node transition, allowing long-running agent flows to resume later without restarting the pipeline.

\begin{center}
\noindent\rule{0.6\textwidth}{0.4pt} \\
\vspace{0.3em}
\textit{[Code implementation is available in the companion coding handbook: \texttt{coding-handbook/ch12\_agentic\_sdks/}]} \\
\noindent\rule{0.6\textwidth}{0.4pt}
\end{center}
